\appendix
\part{Anhang}

\newcommand{\lit}[3]{\vspace{5pt}\noindent\textbf{[#1]} #2\par\nopagebreak\vspace{1pt}{\small\itshape\color{grau} #3}\par\vspace{2pt}}

\chapter{Formelsammlung und Übersicht}

\section{Das Problem}

\begin{merksatz}
\boxtitel{Teil I}
\begin{center}
{\small
\renewcommand{\arraystretch}{1.62}
\begin{tabular}{@{}P{4.4cm}P{4.2cm}@{}}
\toprule
$\Risk(h) = \E[L(h(X),Y)]$ & Risiko (1.2)\\
$h^{*}(x) = [\eta(x) \ge \tfrac12]$ & Bayes-Regel (1.3)\\
$\eta \ge c_{\text{FP}}/(c_{\text{FP}}+c_{\text{FN}})$ & mit Kosten\\
kein Verfahren überall besser & Satz 1.5\\
$\sigma^{2} + \text{Verzerrung}^{2} + \text{Varianz}$ & Satz 2.1\\
$2\sigma^{2}p/n$ & Optimismus (2.2)\\
$c^{*} = \mu^{2}/(\mu^{2}+\sigma^{2})$ & Schrumpfen\\
$e_i/(1-H_{ii})$ & Weglassen (3.2)\\
Anteil richtig geordneter Paare & Trennschärfe (3.3)\\
$\sqrt{2\ln m/n}$ & höchste Zufallskorrelation\\
\bottomrule
\end{tabular}}
\end{center}
\end{merksatz}

\section{Die Verfahren}

\begin{merksatz}
\boxtitel{Teil II und IV}
\begin{center}
{\small
\renewcommand{\arraystretch}{1.62}
\begin{tabular}{@{}P{4.4cm}P{4.2cm}@{}}
\toprule
$1/(1-\rho^{2})$ & Kollinearität\\
$(X^{\!\top}X+\lambda I)^{-1}X^{\!\top}\vek y$ & Kamm (4.2)\\
$d_j^{2}/(d_j^{2}+\lambda)$ & Schrumpffaktor\\
$\operatorname{sign}(b)(|b|-\lambda)_{+}$ & Lasso (4.4)\\
$\ln\frac{p}{1-p} = \beta_0 + \vek\beta^{\!\top}\vek x$ & logistisch (5.1)\\
$\nabla J = \sum(p_i-y_i)\vek x_i$ & Gradient (5.2)\\
$(R/\gamma)^{2}$ & Perzeptron (5.3)\\
$\vek w = \sum\alpha_i t_i\vek x_i$ & Stützvektoren (6.2)\\
$e^{-\gamma\|\vek x-\vek z\|^{2}}$ & Gauß-Kern\\
$-p\log_2p - (1-p)\log_2(1-p)$ & Entropie (9.1)\\
$\rho\sigma^{2} + (1-\rho)\sigma^{2}/B$ & Mittel (9.2)\\
$\prod 2\sqrt{\varepsilon_s(1-\varepsilon_s)}$ & Boosting (9.3)\\
$2\Risk^{*}(1-\Risk^{*})$ & ein Nachbar (10.2)\\
$e = r^{1/p}$ & Kantenlänge\\
$\sigma(1-\sigma) \le \tfrac14$ & verschwindender Gradient\\
\bottomrule
\end{tabular}}
\end{center}
\end{merksatz}

\section{Theorie und Tiefe}

\begin{merksatz}
\boxtitel{Teil III und V}
\begin{center}
{\small
\renewcommand{\arraystretch}{1.62}
\begin{tabular}{@{}P{4.4cm}P{4.2cm}@{}}
\toprule
$(\ln|\Hyp| + \ln\frac1\delta)/\varepsilon$ & fehlerfrei (7.2)\\
$(\ln|\Hyp| + \ln\frac2\delta)/(2\varepsilon^{2})$ & allgemein (7.3)\\
$2e^{-2n\varepsilon^{2}}$ & Hoeffding\\
$\sum_{i\le d}\binom ni \le (en/d)^{d}$ & Sauer (7.5)\\
$\hat\Risk + \sqrt{\frac8n(d\ln\frac{2en}{d}+\ln\frac4\delta)}$ & VC (7.6)\\
Stützvektoren$/n$ & Weglassen (8.1)\\
$(k\ln n + \ln\frac1\delta)/(n-k)$ & Kompression (8.2)\\
$d \le (R/\gamma)^{2}$ & Rand (8.3)\\
$\sum_j\sum_{C_j}\|\vek x_i-\vek\mu_j\|^{2}$ & $k$-Mittel (12.1)\\
$\max\vek u^{\!\top}\Sigma\vek u = \lambda_1$ & Hauptkomp.\ (13.1)\\
$1+\rho(p-1)$ und $1-\rho$ & gleiche Verbundenheit\\
$X^{\!\top}(XX^{\!\top})^{-1}\vek y$ & kürzeste Lösung (14.1)\\
$m+1$ gegen $2^{L}$ Stücke & Tiefe (14.2)\\
\bottomrule
\end{tabular}}
\end{center}
\end{merksatz}

\section{Was in diesem Band \kursiv{nicht} bewiesen wurde}

\begin{merksatz}
\boxtitel{Neun Auslassungen}
\begin{center}
{\small
\renewcommand{\arraystretch}{1.5}
\begin{tabular}{@{}P{3.6cm}P{1.2cm}P{3.2cm}@{}}
\toprule
\textbf{Ergebnis} & \textbf{Kap.} & \textbf{Wo es steht}\\
\midrule
Weglassformel & $3$ & Hastie u.\,a.\zit{11}, Kap.\ 7\\
Satz von Mercer & $6$ & Schölkopf, Smola\zit{22}\\
Hoeffding & $7$ & Shalev-Shwartz\zit{24}, Anh.\ B\\
Sauer & $7$ & Shalev-Shwartz\zit{24}, Kap.\ 6\\
Fundamentalsatz (VC) & $7$ & Shalev-Shwartz\zit{24}, Kap.\ 28\\
Grenzübergang Nachbar & $10$ & Cover, Hart\zit{6}\\
Universelle Näherung & $11$ & Goodfellow u.\,a.\zit{10}\\
EM-Monotonie & $12$ & Bishop\zit{3}, Kap.\ 9\\
Eckart-Young & $13$ & Deisenroth u.\,a.\zit{7}\\
\bottomrule
\end{tabular}}
\end{center}

\begin{center}
\begin{minipage}{0.88\textwidth}
\kursiv{Neun Auslassungen, jede im Text benannt, mit Begründung und Fundstelle. \textbf{Alle übrigen Sätze des Bandes sind vollständig bewiesen} --- darunter die Kompressionsschranke, die Randdimension über das Perzeptron, AdaBoost, die Schranke über Stützvektoren und die kürzeste Lösung des Gradientenabstiegs.}
\end{minipage}
\end{center}
\end{merksatz}

\section{Die Landkarte des Bandes}

\begin{merksatz}
\boxtitel{Welcher Satz worauf beruht}
\begin{center}
{\small
\renewcommand{\arraystretch}{1.4}
\begin{tabular}{@{}P{3.4cm}P{1.0cm}P{3.8cm}@{}}
\toprule
\textbf{Ergebnis} & \textbf{Kap.} & \textbf{beruht auf}\\
\midrule
Bayes-Regel & $1$ & punktweise Minimierung\\
Zerlegung & $2$ & Quadrieren, Erwartungswert\\
Optimismus & $2$ & Spur einer Projektion\\
Kamm & $4$ & positiv definite Matrix\\
Lasso & $4$ & Fallunterscheidung am Knick\\
Konvexität & $5$ & $\sigma' = \sigma(1-\sigma)$\\
Perzeptron & $5$ & Cauchy-Schwarz\\
Stützvektoren & $6$ & komplementärer Schlupf (Band 19)\\
endliche Klassen & $7$ & Vereinigung, $1-\varepsilon \le e^{-\varepsilon}$\\
Weglassen & $8$ & komplementärer Schlupf\\
Kompression & $8$ & wie Satz 7.2\\
Randdimension & $8$ & Perzeptron, Satz 5.3\\
Mittel & $9$ & Kovarianzen summieren\\
Boosting & $9$ & Exponent über Indikator\\
ein Nachbar & $10$ & Varianz $\ge 0$\\
$k$-Mittel & $12$ & Durchschnitt minimiert\\
Hauptkomponenten & $13$ & Spektralsatz (Band 2)\\
kürzeste Lösung & $14$ & Zeilenraum $\perp$ Kern\\
Tiefe & $14$ & Knicke zählen\\
\bottomrule
\end{tabular}}
\end{center}
\end{merksatz}

\section{Zeittafel}

\begin{merksatz}
\boxtitel{Von Hume bis zu den Skalengesetzen}
\begin{center}
{\small
\renewcommand{\arraystretch}{1.32}
\begin{tabular}{@{}P{1.3cm}P{3.0cm}P{4.0cm}@{}}
\toprule
\textbf{Jahr} & \textbf{Person} & \textbf{Ereignis}\\
\midrule
$1739$ & Hume & das Induktionsproblem\\
$1838$ & Verhulst & die logistische Kurve\\
$1894$ & Pearson & Mischung zweier Krabbenarten\\
$1901$ & Pearson & erste Hauptkomponente\\
$1904$ & Spearman & der allgemeine Faktor\\
$1933$ & Hotelling & Hauptkomponenten als Eigenwerte\\
$1936$ & Fisher & Diskriminanzanalyse\\
$1936$ & Eckart, Young & beste Näherung\\
$1943$ & Tichonow & Regularisierung\\
$1950$ & Turing & lernende Maschinen\\
$1951$ & Fix, Hodges & nächste Nachbarn\\
$1956$ & Stein & das Paradoxon\\
$1957$ & Bellman & Fluch der Dimension\\
$1957$ & Lloyd & $k$ Mittelpunkte\\
$1958$ & Rosenblatt & das Perzeptron\\
$1959$ & Samuel & der Name\\
$1962$ & Novikoff & Konvergenz des Perzeptrons\\
$1963$ & Vapnik, Tscherwonenkis & breitester Rand\\
$1967$ & Cover, Hart & ein Nachbar, doppeltes Risiko\\
$1969$ & Minsky, Papert & \kursiv{Perceptrons}\\
$1970$ & Hoerl, Kennard & Kammregression\\
$1971$ & \textbf{Vapnik, Tscherwonenkis} & \textbf{VC-Dimension}\\
$1972$ & Sauer, Shelah & das Lemma\\
$1974$ & Stone & Kreuzvalidierung\\
$1977$ & Dempster, Laird, Rubin & EM\\
$1984$ & Valiant & PAC-Lernen\\
$1984$ & Breiman u.\,a. & Entscheidungsbäume\\
$1986$ & Rumelhart, Hinton, Williams & Rückpropagation\\
$1986$ & Littlestone, Warmuth & Kompression\\
$1989$ & Cybenko & universelle Näherung\\
$1992$ & Boser, Guyon, Vapnik & Kerntrick\\
$1995$ & Freund, Schapire & AdaBoost\\
$1996$ & Tibshirani & Lasso\\
$2001$ & Breiman & Zufallswald, zwei Kulturen\\
$2002$ & Ambroise, McLachlan & das Leck\\
$2012$ & Krizhevsky u.\,a. & tiefe Netze gewinnen\\
$2016$ & Telgarsky & Tiefe\\
$2017$ & \textbf{Zhang u.\,a.} & \textbf{vertauschte Etiketten}\\
$2019$ & Belkin u.\,a. & doppelter Abstieg\\
$2019$ & Obermeyer u.\,a. & Kosten statt Bedarf\\
$2019$ & Seymour u.\,a. & Sepsis-Phänotypen\\
$2020$ & --- & Skalengesetze\\
$2021$ & Wong u.\,a. & Nachprüfung des Sepsissystems\\
$2023$ & Kapoor, Narayanan & Lecks in $17$ Fächern\\
\bottomrule
\end{tabular}}
\end{center}
\end{merksatz}

\section{Glossar}

\begin{merksatz}
\boxtitel{Begriffe von A bis Z}
\begin{center}
{\small
\renewcommand{\arraystretch}{1.3}
\begin{longtable}{@{}P{3.6cm}P{5.0cm}@{}}
\toprule
\textbf{Begriff} & \textbf{Bedeutung}\\
\midrule
\endhead
Äußere Prüfung & an Daten einer anderen Klinik (3)\\
Bagging & Mittel über Neuziehungen (9)\\
Bayes-Regel & überall das Wahrscheinlichere (1.3)\\
Bayes-Risiko & kleinstmöglicher Fehler (1)\\
Doppelter Abstieg & zweites Tal jenseits $p = n$ (14)\\
EM-Verfahren & Erwartung, Maximierung (12)\\
Ersatzverlust & konvexe Näherung an Null-Eins (6)\\
Gauß-Kern & Ähnlichkeit über Abstand (6)\\
Hauptkomponente & Richtung größter Streuung (13)\\
Implizite Regularisierung & Strafe im Verfahren, nicht im Ziel (14)\\
Induktive Neigung & eingebaute Annahme (1)\\
Kalibrierung & $30\,\%$ heißt $30\,\%$ (3)\\
Kammregression & Strafe $\lambda\|\vek\beta\|^{2}$ (4.1)\\
Kern & Skalarprodukt im Merkmalsraum (6.3)\\
Kompression & Regel aus $k$ Punkten (8)\\
Kreuzvalidierung & reihum zurückhalten (3.1)\\
Kreuzentropie & Verlust der log.\ Regression (5)\\
Lasso & Strafe $\lambda\|\vek\beta\|_1$ (4.3)\\
Leck & Information aus der Prüfung im Lernen (3)\\
Nächste Nachbarn & Mehrheit der ähnlichsten (10.1)\\
Optimismus & Test minus Training (2.2)\\
PAC & wahrscheinlich ungefähr richtig (7.1)\\
Perzeptron & Korrektur bei Fehler (5)\\
Rand & Abstand zur Grenze (6.1)\\
Risiko & erwarteter Verlust (1.2)\\
Rückpropagation & Kettenregel rückwärts (11)\\
Stützvektor & Punkt auf dem Rand (6.2)\\
Trennschärfe & Paarvergleich, AUC (3.3)\\
Überanpassung & Varianz überwiegt (2)\\
Unteranpassung & Verzerrung überwiegt (2)\\
VC-Dimension & größte zerschmetterte Menge (7.4)\\
Verzerrung & mittlere Abweichung (2)\\
Zerschmettern & jede Einteilung erreichen (7.4)\\
Zufallswald & Bäume mit Zufallsauswahl (9)\\
\bottomrule
\end{longtable}}
\end{center}
\end{merksatz}
